% !TEX program = lualatex
% Generated from Markdown by src/core/latex/export.ts. Edit the Markdown source, not this file.
\documentclass[12pt,a4paper]{article}
\usepackage{iftex}
\ifPDFTeX
  \usepackage[T1]{fontenc}
  \usepackage[utf8]{inputenc}
  \usepackage{lmodern}
\else
  \usepackage{fontspec}
\fi
\usepackage[english,main=brazilian]{babel}
\usepackage{csquotes}
% ABNT NBR 14724:2024 §5.1: A4, margins 3 cm (top, left) and 2 cm (bottom, right), 1.5 line spacing.
\usepackage[top=3cm,left=3cm,bottom=2cm,right=2cm]{geometry}
\usepackage{setspace}
\onehalfspacing
\usepackage{mathtools}
\usepackage{amssymb}
\usepackage{amsthm}
\usepackage{graphicx}
\usepackage{booktabs}
\usepackage{enumitem}
\usepackage[normalem]{ulem}
\usepackage{fancyvrb}
\usepackage{newfloat}
\usepackage{caption}
% ABNT NBR 14724:2024 §5.8: designation, number, dash and title ABOVE; source below.
\DeclareCaptionLabelSeparator{ft}{ – }
\captionsetup{labelsep=ft,position=top}
\DeclareFloatingEnvironment[fileext=lgr,name={Gráfico},listname={Lista de gráficos}]{ftchart}
\DeclareFloatingEnvironment[fileext=lqd,name={Quadro},listname={Lista de quadros}]{ftframe}
\newcommand{\ftsource}[1]{\par\smallskip{\footnotesize Fonte: #1\par}}
\DefineVerbatimEnvironment{ftverbatim}{Verbatim}{}
\usepackage[backend=biber,style=abnt]{biblatex}
\addbibresource{bibliography.bib}
\usepackage{hyperref}
\hypersetup{hidelinks,pdftitle={Amostragem \& informação: 100\% \textquotedbl{}sinal\textquotedbl{}},pdfauthor={Flavio Tinoco},pdflang={pt-BR}}

% Theorem-like environments: one shared counter (amsthm \newtheorem{x}[theorem]).
\theoremstyle{plain}
\newtheorem{theorem}{Teorema}
\newtheorem{lemma}[theorem]{Lema}
\newtheorem{proposition}[theorem]{Proposição}
\newtheorem{corollary}[theorem]{Corolário}
\newtheorem{conjecture}[theorem]{Conjectura}
\newtheorem{axiom}[theorem]{Axioma}
\theoremstyle{definition}
\newtheorem{definition}[theorem]{Definição}
\newtheorem{example}[theorem]{Exemplo}
\theoremstyle{remark}
\newtheorem*{remark}{Observação}
\renewcommand{\qedsymbol}{\ensuremath{\blacksquare}}

% Markup of the abnt stage (semantic wrappers; see docs/guia-latex.md).
\newcommand{\ftforeign}[2]{\textit{#2}}
\newcommand{\ftname}[1]{#1}
\newcommand{\ftabbr}[2]{#1}
% KaTeX-only HTML extensions degrade to their content.
\providecommand{\htmlClass}[2]{#2}
\providecommand{\htmlData}[2]{#2}
\providecommand{\htmlId}[2]{#2}
\providecommand{\htmlStyle}[2]{#2}

% Site notation (src/core/latex/macros.ts) and document macros.
\makeatletter
\newcommand*{\ftmacro}[3]{%
  \ifdefined#1%
    \expandafter\NewCommandCopy\csname fttext\string#1\endcsname#1%
    \expandafter\newcommand\csname ftmath\string#1\endcsname[#2]{#3}%
    \DeclareRobustCommand{#1}{%
      \ifmmode\expandafter\@firstoftwo\else\expandafter\@secondoftwo\fi
      {\csname ftmath\string#1\endcsname}{\csname fttext\string#1\endcsname}}%
  \else
    \newcommand{#1}[#2]{#3}%
  \fi}
\makeatother
\ftmacro{\N}{0}{\mathbb{N}}
\ftmacro{\Z}{0}{\mathbb{Z}}
\ftmacro{\Q}{0}{\mathbb{Q}}
\ftmacro{\R}{0}{\mathbb{R}}
\ftmacro{\C}{0}{\mathbb{C}}
\ftmacro{\E}{0}{\operatorname{\mathbb{E}}}
\ftmacro{\Var}{0}{\operatorname{Var}}
\ftmacro{\Prob}{0}{\operatorname{P}}
\ftmacro{\H}{0}{\operatorname{H}}
\ftmacro{\I}{0}{\operatorname{I}}
\ftmacro{\KL}{0}{D_{\mathrm{KL}}\infdivx}
\ftmacro{\infdivx}{2}{\left(#1\,\middle\|\,#2\right)}
\ftmacro{\d}{0}{\mathop{}\!\mathrm{d}}
\ftmacro{\e}{0}{\mathrm{e}}
\ftmacro{\i}{0}{\mathrm{i}}
\ftmacro{\sinc}{0}{\operatorname{sinc}}
\ftmacro{\abs}{1}{\left\lvert #1 \right\rvert}
\ftmacro{\norm}{1}{\left\lVert #1 \right\rVert}
\ftmacro{\set}{1}{\left\{ #1 \right\}}
\ftmacro{\F}{0}{\mathcal{F}}
\ftmacro{\defeq}{0}{\coloneqq}
\ftmacro{\entails}{0}{\vDash}
\ftmacro{\proves}{0}{\vdash}

\title{Amostragem \& informação: 100\% \textquotedbl{}sinal\textquotedbl{}\\[0.5ex]\large Um ensaio de teste}
\author{Flavio Tinoco}
\date{14 de março de 2026}

\begin{document}
\maketitle

\begin{abstract}
Resumo com \$ e \#.

\noindent\textbf{Palavras-chave:} amostragem; entropia.
\end{abstract}

\section{Amostragem e informação}\label{amostragem}
% attributes: .claim

A teoria de Shannon \parencite[{p.~379}]{shannon1948} mede a informação em bits; ver também
\parencites[{ver}][{}]{nyquist1928}[{cap.~2}]{kotelnikov1933} e a releitura de \textcite[{p.~12}]{gallager2008}.
Já \parencite*{shannon1949} tratou do ruído. Escrever \texttt{R\$ 5,00 \& 10\%} não é fórmula, e
\textquotedbl{}aspas\textquotedbl{} com \# \$ \% \& \textasciitilde{} \_ \textasciicircum{} \textbackslash{} \{ \} devem sobreviver; o e-mail \href{mailto:fulano@exemplo.com}{fulano@exemplo.com} também.
Às 10:30 a proporção é 1:2 e [@escapado] não é citação.

\begin{definition}[{Entropia}]\label{def:entropia}
A \textbf{entropia} de uma variável $X$ é

\begin{equation}
\H(X) = -\sum_{x} p(x) \log_2 p(x) \label{eq:entropia}
\end{equation}
\end{definition}

\begin{theorem}[{Nyquist–Shannon}]\label{thm:amostragem}
Se $x$ é limitado em banda a $B$, então $x$ é determinado por amostras a $2B$
(compare \eqref{eq:entropia} e Definição~\ref{def:entropia}).
\end{theorem}

\begin{proof}[{Demonstração do Teorema~\ref{thm:amostragem}}]
Pela Equação~\eqref{eq:sinc}, segue o resultado.
\end{proof}

\begin{equation}
x(t) = \sum_{n} x(nT) \sinc\left(\frac{t - nT}{T}\right) \label{eq:sinc}
\end{equation}

\begin{ftchart}[htbp]
\centering
\caption{Dobramento espectral}\label{fig:dobramento}
\includegraphics[width=\linewidth,alt={Espectro dobrado em torno de f\_s/2}]{./dobramento}
\ftsource{Elaborado pelo autor (2026).}
\end{ftchart}

\begin{table}[htbp]
\centering
\caption{Taxas de amostragem usuais}\label{tab:taxas}
\begin{tabular}{@{}lcr@{}}
\toprule
Sinal & Banda (kHz) & Taxa (kHz) \\
\midrule
Voz & 4 & 8 \\
CD & 22,05 & 44,1 \\
\bottomrule
\end{tabular}
\ftsource{\parencite{iso80000}}
\end{table}

\begin{remark}
Uma \emph{observação} com\footnote{nota lateral} e \ftforeign{de}{Zeitgeist},
\ftname{\mbox{Ferdinand} \mbox{de} \mbox{Saussure}} e \ftabbr{FFT}{Fast Fourier Transform}.
\end{remark}

\begin{enumerate}
\item primeiro item
\item {}[segundo] item com colchete
\end{enumerate}

\begin{itemize}
\item[$\boxtimes$] tarefa feita
\item[$\square$] tarefa pendente
\end{itemize}

\begin{quote}
Citação em bloco.\footnote{Nota de rodapé com \texttt{código}.}
\end{quote}

% language: ts
\begin{ftverbatim}
const x = 1; // \end{verbatim}
\end{ftverbatim}

\tableofcontents

\printbibliography

\end{document}
